% Volume V: Dynamics, Signals, Graphs, and Control
% Continuity note for Chapter 30.
% Previous dependencies: Chapters 25 and 27 provide deterministic and stochastic dynamics; Chapter 29 provides state-transition graph intuition; Chapters 15 and 17 provide conditioning, Bayes decisions, losses, and risks.
% New durable concepts: controlled state-space systems; state, control/action, observation, feedback, policy, plant, controllability intuition, cost functionals, running and terminal costs, optimal control, minimum-energy control, LQR intuition, finite-horizon dynamic programming, Bellman equations, infinite-horizon discounting, value functions, Q-functions, reinforcement-learning links, biological feedback, delays, noise, internal models, and neural decision dynamics.
% Forward hooks: Chapter 36 builds reinforcement learning from Markov decision processes and Bellman equations; Chapter 45 relates plasticity to learning rules; Chapter 50 connects decision-making, Bayesian inference, and predictive coding.
% Notation commitments: Continuous controlled states use \dot x=f(x,u,t); linear systems use \dot x=Ax+Bu; discrete states use s_t and actions a_t; policies are \pi(a\mid s) or u_t=\pi_t(s_t); value functions are V_t(s) and Q_t(s,a).
\section{Chapter 30: Control Theory and Dynamic Programming}
\label{ch:control-theory-dynamic-programming}

Dynamics describe how a state changes. Control asks how to choose inputs so that the state changes in a desired way. This is the mathematical language of reaching movements, eye stabilization, homeostasis, robotics, optimal feedback, and reinforcement learning. It also clarifies a basic modeling distinction: predicting a system and controlling a system are related but not identical tasks.

\paragraph{Chapter problem.}
How do we represent actions in dynamical systems, define objectives over trajectories, compute optimal decisions recursively, and interpret control ideas in biological and artificial agents?

\paragraph{Section map.}
Section 30.1 defines controlled dynamical systems, feedback, linear control, and controllability intuition. Section 30.2 formulates optimal control as minimizing trajectory cost and works through a minimum-energy mechanism. Section 30.3 derives Bellman equations for finite-horizon and infinite-horizon dynamic programming and connects them to reinforcement learning. Section 30.4 applies feedback, delay, noise, internal models, and decision dynamics to biological systems without overclaiming the implementation.

\subsection{30.1 Controlled dynamical systems}

\paragraph{Guiding question.}
How does a dynamical-system model change when an agent or controller can choose inputs?

\paragraph{Beginner-facing intuition.}
An uncontrolled system evolves according to its own dynamics. A controlled system has inputs that can push it. The state might be hand position, body temperature, membrane voltage, robot configuration, or a latent neural state. The control might be muscle activation, heater power, injected current, torque, or an action chosen by a policy.

The key distinction is between an open-loop input, chosen as a function of time only, and feedback, chosen using the current or estimated state. Feedback is powerful because it can correct errors caused by noise, perturbations, or wrong initial conditions.

\paragraph{Formal definitions and notation.}
A continuous-time controlled system has state
\[
x(t)\in\mathbb R^d
\]
and control input
\[
u(t)\in\mathbb R^m.
\]
Its dynamics are
\[
\boxed{\dot x(t)=f(x(t),u(t),t).}
\]
A linear time-invariant controlled system has
\[
\boxed{\dot x(t)=Ax(t)+Bu(t),}
\]
where \(A\in\mathbb R^{d\times d}\) describes intrinsic dynamics and \(B\in\mathbb R^{d\times m}\) describes how controls enter the state. A feedback controller has the form
\[
u(t)=\pi(x(t),t)
\]
or, in a linear state-feedback law,
\[
u(t)=Kx(t).
\]
If the state is not directly observed, a controller may use an estimate \(\widehat x(t)\), connecting this chapter to filtering in Chapter 27.

\paragraph{Core mechanism: feedback changes eigenvalues.}
For the linear system
\[
\dot x=Ax+Bu
\]
with feedback \(u=Kx\), the closed-loop system is
\[
\dot x=(A+BK)x.
\]
Thus the feedback gain changes the effective dynamics. In one dimension,
\[
\dot x=ax+bu,\qquad u=kx
\]
becomes
\[
\dot x=(a+bk)x.
\]
The equilibrium \(x=0\) is stable when \(a+bk<0\). If \(b\neq0\), an appropriate gain \(k\) can move the scalar closed-loop eigenvalue to the left half-line. In higher dimensions, controllability determines which modes can be moved by feedback.

\paragraph{Concrete example.}
Let
\[
\dot x=x+u.
\]
Without control, \(x(t)=x(0)e^t\) grows away from \(0\). With feedback \(u=-2x\),
\[
\dot x=-x,
\]
so \(x(t)=x(0)e^{-t}\) decays toward \(0\). The same physical state variable has different dynamics because the controller closes a feedback loop.

\paragraph{Computational neuroscience example.}
In a reaching movement, the state may include hand position and velocity, while the control is muscle activation. Sensory feedback can correct deviations caused by perturbations. A model
\[
\dot x=Ax+Bu+\text{noise}
\]
does not claim the nervous system explicitly stores matrices \(A\) and \(B\); it provides a precise way to separate plant dynamics, control inputs, noise, and feedback.

\paragraph{AI or machine-learning example.}
In robotics, the policy \(u=\pi(x)\) maps measured or estimated state to torque commands. In model-based reinforcement learning, a learned dynamics model predicts \(x_{t+1}\) from \(x_t\) and action \(u_t\), and a planner chooses controls. In recurrent neural networks, external inputs can be interpreted as controls that steer hidden states.

\paragraph{Common misconceptions and failure modes.}
Control is not just forcing a desired trajectory; constraints, noise, and input limits matter. Feedback can destabilize a system if the gain or delay is wrong. Observing a correlation between action and state change does not identify \(B\) without an intervention or model. Stabilizability of a linear system depends on whether unstable modes are reachable through \(B\).

\paragraph{Exercises.}
\begin{enumerate}[label=\textbf{30.1.\arabic*.}, leftmargin=*]
\item \textbf{Mechanical.} In \(\dot x=Ax+Bu\), state the shapes of \(A\) and \(B\) when \(x\in\mathbb R^4\) and \(u\in\mathbb R^2\).
\item \textbf{Proof or derivation.} Derive the closed-loop matrix \(A+BK\) when \(u=Kx\).
\item \textbf{Numerical/computational.} For \(\dot x=3x+2u\), choose a scalar feedback \(u=kx\) that makes the closed-loop coefficient equal to \(-1\).
\item \textbf{Interpretation.} Why is feedback useful when the initial state or perturbations are uncertain?
\item \textbf{Misconception/debugging.} A student says, ``Large feedback gains always improve control.'' Explain how delay, noise amplification, or instability can make this false.
\end{enumerate}

\subsection{30.2 Optimal control}

\paragraph{Guiding question.}
How do we choose controls that balance task success against cost, effort, time, and risk?

\paragraph{Beginner-facing intuition.}
Many control problems have more than one way to reach the goal. A fast movement may require large effort. A cautious policy may be accurate but slow. Optimal control makes this tradeoff explicit by assigning a cost to trajectories and controls, then choosing the control that minimizes expected or deterministic cost.

\paragraph{Formal definitions and notation.}
For continuous-time dynamics
\[
\dot x(t)=f(x(t),u(t),t),\qquad x(0)=x_0,
\]
an optimal-control objective often has the form
\[
\boxed{
J(u)=\int_0^T c(x(t),u(t),t)\,dt+\phi(x(T)).
}
\]
The function \(c\) is a running cost and \(\phi\) is a terminal cost. The goal is
\[
u^\star=\arg\min_{u(\cdot)}J(u)
\]
subject to the dynamics and any control constraints. In stochastic control, \(J\) is usually an expected cost.

The linear-quadratic regulator, or LQR, is the canonical case:
\[
\dot x=Ax+Bu,\qquad
c(x,u)=x^\top Qx+u^\top Ru,
\]
with \(Q\succeq0\) and \(R\succ0\). Its optimal feedback is linear under standard assumptions.

\paragraph{Core mechanism: minimum-energy control for an integrator.}
Consider the scalar integrator
\[
\dot x(t)=u(t),\qquad x(0)=x_0,\qquad x(T)=x_T.
\]
Minimize energy
\[
J(u)=\int_0^T u(t)^2\,dt.
\]
The constraint implies
\[
x_T-x_0=\int_0^T u(t)\,dt.
\]
By Cauchy--Schwarz,
\[
\left(\int_0^T u(t)\,dt\right)^2
\leq
\left(\int_0^T 1^2\,dt\right)
\left(\int_0^T u(t)^2\,dt\right)
=TJ(u).
\]
Therefore
\[
J(u)\geq \frac{(x_T-x_0)^2}{T}.
\]
Equality holds when \(u(t)\) is constant:
\[
u^\star(t)=\frac{x_T-x_0}{T}.
\]
The least-energy way to move an integrator between fixed endpoints in fixed time is to spread effort evenly.

\paragraph{Concrete example.}
Move from \(x_0=0\) to \(x_T=10\) in \(T=5\) s with \(\dot x=u\). The minimum-energy control is
\[
u^\star(t)=2.
\]
The minimum cost is
\[
J^\star=\int_0^5 4\,dt=20.
\]
Using \(u=10\) for one second and \(u=0\) afterward reaches the same endpoint, but costs \(100\), because concentrated effort is expensive under a quadratic cost.

\paragraph{Computational neuroscience example.}
Optimal feedback control models reaching by balancing endpoint accuracy, movement duration, signal-dependent noise, and control effort. The mathematics does not imply the motor system explicitly solves the exact LQR equations on every movement. It states a normative model: if a system minimizes this cost under these dynamics and noise assumptions, then the optimal feedback has a particular structure that can be compared with behavior.

\paragraph{AI or machine-learning example.}
Model predictive control repeatedly solves a finite-horizon optimal-control problem using the current state estimate, applies the first action, then replans. This approach is used in robotics and can be combined with learned dynamics models. The cost function encodes the task: tracking, safety, energy use, or constraint violation.

\paragraph{Common misconceptions and failure modes.}
Optimal means optimal for the specified model and cost, not universally best. A cost function is a modeling choice. If effort cost is too high, the controller may prefer inaction. If terminal cost is too low, it may miss the target. If the dynamics model is wrong, the computed control can perform poorly. Constraints can change the solution qualitatively.

\paragraph{Exercises.}
\begin{enumerate}[label=\textbf{30.2.\arabic*.}, leftmargin=*]
\item \textbf{Mechanical.} In \(J(u)=\int_0^T c(x,u,t)\,dt+\phi(x(T))\), identify the running cost and terminal cost.
\item \textbf{Proof or derivation.} Reproduce the Cauchy--Schwarz derivation for minimum-energy control of \(\dot x=u\).
\item \textbf{Numerical/computational.} Move from \(x_0=1\) to \(x_T=7\) in \(T=3\) under \(\dot x=u\). Compute \(u^\star\) and \(J^\star\).
\item \textbf{Interpretation.} Why does a quadratic effort cost favor spreading control over time rather than using a brief large pulse?
\item \textbf{Misconception/debugging.} A student fits a behavior with an optimal-control model and says, ``The brain must minimize exactly this cost.'' Explain the difference between a normative model and a mechanistic claim.
\end{enumerate}

\subsection{30.3 Dynamic programming}

\paragraph{Guiding question.}
How can a long decision problem be solved by working backward from future value?

\paragraph{Beginner-facing intuition.}
Dynamic programming is the mathematics of reusable subproblems. If the future consequences of an action depend only on the next state, then the value of being in a state summarizes everything needed for optimal future choice. This is the principle of optimality: an optimal plan has optimal continuation plans.

This is the core bridge from control theory to reinforcement learning. Bellman equations are not just formulas; they are recursive definitions of optimal future cost or reward.

\paragraph{Formal definitions and notation.}
A finite-horizon Markov decision process has states \(s\in S\), actions \(a\in A(s)\), transition probabilities
\[
p(s'\mid s,a),
\]
one-step costs \(c_t(s,a)\), and terminal cost \(g(s)\). A deterministic Markov policy chooses
\[
a_t=\pi_t(s_t).
\]
The optimal finite-horizon value function is
\[
V_t(s)=\min_{\pi}\mathbb E\left[\sum_{k=t}^{T-1}c_k(S_k,A_k)+g(S_T)\mid S_t=s\right].
\]
The Bellman recursion is
\[
\boxed{
V_T(s)=g(s),
\qquad
V_t(s)=\min_{a\in A(s)}
\left[
c_t(s,a)+\sum_{s'}p(s'\mid s,a)V_{t+1}(s')
\right].
}
\]
For reward maximization, replace costs by rewards and \(\min\) by \(\max\).

For an infinite-horizon discounted problem with discount \(0<\gamma<1\), the optimal value satisfies
\[
\boxed{
V^\star(s)=\min_a\left[c(s,a)+\gamma\sum_{s'}p(s'\mid s,a)V^\star(s')\right].
}
\]
The state-action value is
\[
Q^\star(s,a)=c(s,a)+\gamma\sum_{s'}p(s'\mid s,a)V^\star(s').
\]

\paragraph{Core mechanism: one Bellman backup.}
Suppose a state \(s\) has two actions. Action \(a_1\) costs \(1\) now and moves to states \(u,v\) with probabilities \(0.5,0.5\). Action \(a_2\) costs \(3\) now and moves to \(u\) with probability \(1\). If
\[
V_{t+1}(u)=4,\qquad V_{t+1}(v)=10,
\]
then
\[
Q_t(s,a_1)=1+0.5(4)+0.5(10)=8,
\]
while
\[
Q_t(s,a_2)=3+4=7.
\]
The Bellman backup chooses \(a_2\), even though its immediate cost is larger, because it leads to better future value.

\paragraph{Concrete example.}
Consider a two-state discounted problem with costs. From state \(0\), action stay costs \(1\) and remains in \(0\); action switch costs \(3\) and moves to \(1\). State \(1\) has zero cost and remains in \(1\). With \(\gamma=0.9\),
\[
V^\star(1)=0.
\]
At state \(0\),
\[
V^\star(0)=\min\{1+0.9V^\star(0),\ 3+0.9V^\star(1)\}
=\min\{1+0.9V^\star(0),3\}.
\]
Always staying would cost \(10\). Switching costs \(3\), so \(V^\star(0)=3\) and the optimal action is switch.

\paragraph{Computational neuroscience example.}
Sequential decision tasks can be modeled with latent states, actions, rewards or costs, and transition probabilities. Neural activity in frontal cortex, basal ganglia, or parietal cortex may correlate with value-like variables in such tasks. The Bellman equation gives a precise computational quantity to compare against neural data, but a correlation with \(V\) or prediction error does not by itself identify the biological circuit implementation.

\paragraph{AI or machine-learning example.}
Reinforcement learning algorithms estimate Bellman quantities from data. Value iteration applies Bellman backups when the model is known. Q-learning estimates \(Q^\star(s,a)\) from sampled transitions. Policy-gradient methods optimize parameterized policies more directly, but still evaluate actions by future return. Chapter 36 develops these ideas after this mathematical foundation.

\paragraph{Common misconceptions and failure modes.}
Dynamic programming requires a state representation that makes the future conditionally independent of the past given the current state and action. If important hidden variables are omitted, the Bellman recursion is applied to the wrong state. Discounting is not merely a numerical trick; it changes the objective by valuing near-term costs or rewards more. Greedy minimization of immediate cost can be suboptimal when future value matters.

\paragraph{Exercises.}
\begin{enumerate}[label=\textbf{30.3.\arabic*.}, leftmargin=*]
\item \textbf{Mechanical.} Define \(V_t(s)\), \(p(s'\mid s,a)\), and \(c_t(s,a)\) in a finite-horizon Markov decision process.
\item \textbf{Proof or derivation.} Derive the finite-horizon Bellman recursion by conditioning on the first action and next state.
\item \textbf{Numerical/computational.} In the one-backup example, recompute the optimal action if \(V_{t+1}(v)=2\) instead of \(10\).
\item \textbf{Interpretation.} Why can an action with higher immediate cost be optimal?
\item \textbf{Misconception/debugging.} A student applies a Bellman equation to a task where the state omits the animal's hidden motivation. What Markov assumption may fail?
\end{enumerate}

\subsection{30.4 Control in biological systems}

\paragraph{Guiding question.}
How can feedback-control and dynamic-programming ideas help model biological behavior without pretending the equations are literal neural code?

\paragraph{Beginner-facing intuition.}
Biological control systems face delays, sensory noise, motor noise, changing bodies, and uncertain goals. They often combine feedback with prediction. A reflex uses fast sensory feedback. An internal model predicts the sensory consequences of action. A decision circuit accumulates noisy evidence before committing to an action. These ideas can be written mathematically without claiming that biology implements the exact algorithms used in engineering.

\paragraph{Formal definitions and notation.}
A feedback loop can be described by a plant state \(x_t\), sensory observation
\[
y_t=h(x_t)+\eta_t,
\]
state estimate \(\widehat x_t\), controller
\[
u_t=\pi(\widehat x_t,t),
\]
and plant dynamics
\[
x_{t+1}=F(x_t,u_t)+\epsilon_t.
\]
Noise variables \(\eta_t\) and \(\epsilon_t\) represent measurement and process uncertainty. Delays mean the controller may use \(y_{t-\delta}\) rather than \(y_t\). A biological cost may include task error, effort, variability, risk, or time.

In decision dynamics, an evidence variable may follow the drift-diffusion model
\[
dZ_t=\mu\,dt+\sigma\,dW_t,
\]
with a choice made when \(Z_t\) reaches an upper or lower threshold. This is a stochastic control-adjacent model of speed-accuracy tradeoffs, not a universal model of all decisions.

\paragraph{Core mechanism: feedback error correction with stability limits.}
Consider a scalar error \(e_t\) updated by proportional feedback:
\[
e_{t+1}=e_t-u_t+\eta_t,\qquad u_t=ke_t.
\]
Then
\[
e_{t+1}=(1-k)e_t+\eta_t.
\]
Ignoring noise, errors shrink when
\[
|1-k|<1,
\]
or
\[
0<k<2.
\]
If \(k\) is too small, correction is slow. If \(k\) is too large, errors alternate signs and grow. With sensory delay, the stable range can be smaller. This simple calculation explains why biological feedback must balance responsiveness against overshoot and noise amplification.

\paragraph{Concrete example.}
Let \(e_0=10\), \(k=0.5\), and \(\eta_t=0\). Then
\[
e_1=5,\qquad e_2=2.5,\qquad e_3=1.25.
\]
With \(k=1.5\),
\[
e_1=-5,\qquad e_2=2.5,\qquad e_3=-1.25.
\]
The second controller converges in magnitude but overshoots each step. With \(k=2.5\), the factor is \(-1.5\), so the error magnitude grows.

\paragraph{Computational neuroscience example.}
Postural control combines delayed sensory feedback, body mechanics, muscle activation, and prediction. Reaching movements combine feedforward plans with feedback corrections. Neural decision tasks often show activity that resembles accumulation to a bound. These are mathematically distinct control motifs: feedback stabilization, optimal movement, and sequential evidence accumulation.

\paragraph{AI or machine-learning example.}
Robotics uses similar motifs: state estimation from sensors, feedback control for stabilization, model predictive control for constrained movement, and reinforcement learning when dynamics or objectives must be learned from interaction. The same Bellman equations can describe artificial agents and normative models of behavior, but the implementation details differ.

\paragraph{Common misconceptions and failure modes.}
Not every negative-feedback loop is optimal. Not every neural correlate of reward prediction error proves model-free reinforcement learning. A good engineering controller may be biologically implausible because it assumes precise state access, no delay, or unlimited control. Conversely, biological systems may use robust heuristics rather than exact solutions to an optimal-control problem.

\paragraph{Exercises.}
\begin{enumerate}[label=\textbf{30.4.\arabic*.}, leftmargin=*]
\item \textbf{Mechanical.} In the feedback loop \(y_t=h(x_t)+\eta_t\), \(u_t=\pi(\widehat x_t,t)\), identify the plant state, observation, estimator, and controller.
\item \textbf{Proof or derivation.} For \(e_{t+1}=(1-k)e_t\), derive the condition on \(k\) for convergence to zero.
\item \textbf{Numerical/computational.} Starting with \(e_0=8\), compute \(e_1,e_2,e_3\) for \(k=0.25\) and no noise.
\item \textbf{Interpretation.} Why do sensory delays make high-gain feedback risky in biological motor control?
\item \textbf{Misconception/debugging.} A student says, ``If behavior matches an optimal-control prediction, the brain must be solving the exact optimization problem.'' State a more careful interpretation.
\end{enumerate}

\paragraph{Closing bridge.}
Control theory turns dynamics into purposeful action by adding inputs, policies, costs, and feedback. Dynamic programming shows how future consequences can be folded backward into present value. Chapter 36 will return to these Bellman ideas in reinforcement learning, where transition models and value functions are learned from data rather than given in advance.

\clearpage
